\scanpage{135}
\mspage{67}

On va donc admettre que (37) définit bien une structure d'\unsure{extensions|extension}, et on aura un hom de suites exactes

\deleted{\unsure{(38)}}
\[
\tag{38}
\diagram{
\begin{tikzcd}[row sep=large,column sep=large]
1 \arrow[r,dash] & T^J \arrow[r] \arrow[d] & \mathcal E_{U_{J,\alpha}}^{\bar L} \arrow[r] \arrow[d] & \mathcal E_{D_{J,\alpha}}^{\tilde D_{J,\alpha}} \arrow[r] \arrow[d] & 1 \\
1 \arrow[r,dash] & \pi_{U/K}^{\bar L} \arrow[r] & \mathcal E_U^{\bar L} \arrow[r] & \mathcal E_K^{\bar K} \arrow[r] & 1
\end{tikzcd}
}
\]
--- où on a rectifié le tir, en faisant attention que (37) n'a de sens telle quelle que si $D_J$ est connexe, \inserted{i.e. $X_J$ connexe} que sinon il faut regarder les composantes connexes $X_{J,\alpha}$ de $X_J$, induisant les composantes \unsure{conn|conn.} $D_{J,\alpha}$ de $D_J$, et $U_{J,\alpha}$ de $U_J$, et écrire la suite exacte pour chaque composante. Bien sûr, il y a des choses plus intrinsèques : écrire avec les topos \deleted{étales \unsure{galoisiens|géométriques}}\inserted{multigaloisiens}, et leurs groupoïdes fondamentaux....\deleted{\unsure{classes de} $\mathcal E_U^{\bar L}$ \unsure{conjugués|conjuguées}}

On voit ainsi que les hom
\[
\tag{39}
T^J\longrightarrow\unsure{\pi|\mathcal E}_{U/K}^{\bar L}
\]
induits par les différents $x\in D_J$ --- \unsure{chacun} n'étant en fait, quand on reprend la définition précise, que défini modulo opérations \unsure{directes|induites} par autom. intérieurs de $\mathcal E_U^{\bar L}$, induits par

\margin{134}
